
\section{Compliance with Standards and Professional Practice}
The project follows IEEE-style software engineering documentation, uses open-source licenses compatible with research distribution, and addresses data privacy ethically by minimizing exposure of exam content. ISO/IEC 25010 qualities (security, maintainability, usability) informed non-functional requirements (NFR1--NFR5).

\subsection{Software Standard}
Python PEP~8 conventions, modular packages (\texttt{ingestion}, \texttt{retriever}, \texttt{privacy}, \texttt{federated}), version-pinned \texttt{requirements.txt}.

\subsection{Hardware Standard}
Consumer PC/lab hardware with CPU inference; no custom PCB design.

\section{Design Constraints}
\subsection{Ethical and Professional Responsibilities}
Results are reported from scripted benchmarks without fabrication. Protected exam snippets in privacy tests are synthetic placeholders. Users must not deploy the demo against real exams without institutional approval.

\subsection{Environmental and Sustainability Considerations}
Quantized SLMs and CPU inference reduce energy versus always-on cloud GPU APIs.

\subsection{Health and Safety Considerations}
Not applicable beyond standard computer-lab safety; no clinical or physical control systems.

\section{Complex Engineering Problem (CEP) Coverage}
\begin{table}[H]
\centering
\caption{CEP mapping summary}
\label{tab:cep_mapping}
\renewcommand{\arraystretch}{1.2}
\small
\begin{tabular}{|c|p{9.5cm}|}
\hline
\textbf{P} & \textbf{Evidence in this project} \\
\hline
P1 & Deep analysis of privacy--utility trade-offs in RAG; Bloom ordinal metrics. \\
\hline
P2 & Conflicting requirements: helpful answers vs exam non-leakage. \\
\hline
P3 & Multi-disciplinary: NLP, IR, security, education theory (Bloom). \\
\hline
P4 & Novel integration of dual federated layers + dual-role SLM stack. \\
\hline
P5 & In-depth evaluation suite and ablation-ready scripts. \\
\hline
P6 & Student/teacher/IT admin stakeholder needs in requirements. \\
\hline
\end{tabular}
\end{table}

\section{Complex Engineering Activity (CEA) Coverage}
\begin{table}[H]
\centering
\caption{CEA mapping summary}
\label{tab:cea_mapping}
\renewcommand{\arraystretch}{1.2}
\small
\begin{tabular}{|c|p{9.5cm}|}
\hline
\textbf{A} & \textbf{Evidence in this project} \\
\hline
A1 & Application of mathematics, CS, and ML theory (embeddings, softmax classification, FedAvg). \\
\hline
A2 & Problem analysis: role separation, adversarial taxonomy, ordinal Bloom metrics. \\
\hline
A3 & Design of modules: ingestion, retriever, PrivacyGuard, federated LoRA/privacy layers. \\
\hline
A4 & Investigation through scripted benchmarks (\texttt{run\_evaluation\_pipeline.py}). \\
\hline
A5 & Modern tool usage: PyMuPDF, Tesseract, FAISS, Transformers, PEFT, Streamlit. \\
\hline
A6 & Professional responsibility: synthetic exam data, ethical deployment notes, reproducible artifacts. \\
\hline
\end{tabular}
\end{table}

\section{Summary}
The project addresses a complex engineering problem in educational AI with traceable CEP/CEA alignment and explicit ethical constraints.
